% !TeX root = ../main.tex
\section{Methodology}
\label{sec:methodology}

The methodology combines a fixed hybrid CIM mapping with controlled perturbation models, so that algorithmic training effects can be separated from device-level precision, variability, and drift constraints.

\subsection{Hybrid Analog/Digital Inference Stack}
\label{subsec:hybrid-deployment}

We adopt a hybrid analog/digital inference stack: dense linear operators execute on differential-pair crossbars; control-heavy operators remain digital. Concretely, the analog partition contains patch embeddings, query/key/value (Q/K/V) and output projections, and feed-forward layers; the digital partition contains softmax, layer normalization, GELU activation, positional encoding, residual additions, and the class token, following established practice in heterogeneous CIM accelerators \citep{wu2023bwq,wang2024epim,ge2024allspark,wang2025hemlet,liu2024hardsea}.

Under this mapping, approximately 88\% of parameters execute on analog arrays while roughly 60\% of estimated energy remains digital, largely because attention operations stay in the digital domain---a pattern consistent with recent heterogeneous CIM accelerators for Vision Transformers \citep{liu2024hardsea,wang2024epim,ge2024allspark,wang2025hemlet}.

\begin{table}[htbp]
\centering
\caption{\textbf{Canonical Simulation Parameters and Hardware Constraints.} These values define the default non-ideality regime for the Tiny-ViT-5M experiments unless otherwise specified.}
\label{tab:hardware_constraints}
\begin{tabular}{lll}
\toprule
\textbf{Parameter} & \textbf{Value} & \textbf{Physical Meaning} \\
\midrule
$n_{\mathrm{states}}$ & 16 & Conductance resolution (4-bit equivalent) \\
$\sigma_{\mathrm{D2D}}$ & 10\% & Device-to-device mismatch (relative std. dev.) \\
$\sigma_{\mathrm{C2C}}$ & 5\% & Cycle-to-cycle read noise \\
$G_{\max}/G_{\min}$ & 10 & Dynamic range of analog weights \\
ADC bits & 8-bit & Default readout quantization (before sweep) \\
\bottomrule
\end{tabular}
\end{table}

\subsection{Modeling Physical Non-Idealities}
\label{subsec:modeling-nonidealities}

The framework injects quantization, cycle-to-cycle (C2C) and device-to-device (D2D) variability, ADC distortion, retention drift, and nonlinear-write surrogates directly into the forward path. Each analog-mapped weight is decomposed into positive and negative branches, mapped to the conductance window $[G_{\min}, G_{\max}]$, and quantized to $n_{\mathrm{states}}$ levels via a straight-through estimator (STE) quantizer $Q_{n}(\cdot)$. Differential readout recovers $G_{\mathrm{eff}}=\tilde{G}^{+}-\tilde{G}^{-}$, and the downstream digital weight is
\begin{equation}
    \hat{W}_{ij}=s_{\ell}\,G_{\mathrm{eff},ij},
    \qquad
    s_{\ell}=
    \begin{cases}
        \dfrac{w_{\max}}{G_{\max}-G_{\min}}, & \text{standard calibration} \\[6pt]
        \dfrac{w_{\max}}{\max_{i,j}|G_{\mathrm{eff},ij}|+\epsilon}, & \text{retention-time recalibration}
    \end{cases}
\label{eq:scale-recovery}
\end{equation}
with layer index $\ell$ omitted. Unless otherwise stated, we use the standard-calibration branch; the retention-recalibration branch is used only for specific retention-drift diagnostics. The recovered weight step is $\Delta_{\hat{W}}=s_{\ell}(G_{\max}-G_{\min})/(n_{\mathrm{states}}-1)$, so perturbations smaller than half a step are masked. The differential-pair scheme suppresses common-mode noise \citep{wei2020voltagedifferential,kim2024sttmram}.

Hardware-aware training (HAT) injects forward perturbations while back-propagating through a straight-through estimator. Standard HAT fixes one D2D mismatch field $M_{0}$ during training.
\paragraph{Ensemble HAT as distribution matching.}
Ensemble HAT treats the mismatch map as a random variable resampled each epoch:
\begin{equation}
    \theta_{\text{ens}}^{\star} = \arg\min_{\theta}\;\mathbb{E}_{M\sim\mathcal{D}_{\text{D2D}}}\Bigl[\frac{1}{N}\sum_{n=1}^{N}\ell\bigl(f(x_n;\theta,M),y_n\bigr)\Bigr],
    \label{eq:hat-ensemble-distribution}
\end{equation}
where the expectation is Monte-Carlo estimated with one fresh draw $M^{(e)}\sim\mathcal{D}_{\text{D2D}}$ held constant across all mini-batches in epoch $e$. Expanding Eq.~\ref{eq:hat-ensemble-distribution} (Supplementary Note S-Theory) reveals an implicit Fisher-weighted gradient-squared regularization effect that pushes the optimizer toward flat regions in the $M$-direction, explaining the fresh-instance transferability without extra hyperparameters. Epoch-level (not per-batch) resampling is load-bearing, giving the optimizer enough consistent signal to fit each instance before resampling \citep{tobin2017domain}. Fresh-instance transfer is then evaluated on unseen i.i.d. D2D draws held constant for each evaluation pass. For a single checkpoint, the reported fresh-instance spread is the standard deviation across ten fresh-instance means (each averaged over five forward passes); for the primary results, the uncertainty is the sample standard deviation across three seed-level fresh-instance means\footnote{The canonical single-seed headline used in the thesis is 86.37 $\pm$ 1.54\%; this study reports the 3-seed mean 86.16 $\pm$ 0.19\% to emphasize reproducibility robustness.}.

\begin{algorithm}[H]
\caption{Ensemble Hardware-Aware Training (HAT)}
\label{alg:ensemble-hat}
\begin{algorithmic}[1]
\Procedure{Train}{$\mathcal{D}_{\text{train}}, \eta, \mathcal{D}_{\text{D2D}}$}
    \State Initialize weights $\theta$
    \For{epoch $e = 1 \dots E$}
        \State Sample fresh mismatch mask $M^{(e)} \sim \mathcal{D}_{\text{D2D}}$
        \For{batch $(x, y) \in \mathcal{D}_{\text{train}}$}
            \State Compute output with current mask: $\hat{y} = f(x; \theta, M^{(e)})$
            \State Update $\theta \leftarrow \theta - \eta \nabla_{\theta} \ell(\hat{y}, y)$
        \EndFor
    \EndFor
    \State \Return $\theta$
\EndProcedure
\Procedure{Evaluate}{$x, \theta, \mathcal{D}_{\text{D2D}}$}
    \State Sample unseen masks $\{M'_{k}\}_{k=1}^K \sim \mathcal{D}_{\text{D2D}}$
    \State \Return $\frac{1}{K} \sum_{k=1}^K \text{Acc}(f(x; \theta, M'_k))$
\EndProcedure
\end{algorithmic}
\end{algorithm}

\subsection{Training Protocol}
\label{subsec:training-protocol}

All Tiny-ViT experiments use AdamW ($lr = 5\times10^{-4}$, weight decay $0.05$), 100 epochs, batch size 64, cosine annealing ($T_{\max}=100$). Images are resized to $224\times224$ with standard augmentation. The default analog resolution is $n_{\mathrm{states}}=16$ levels. We evaluate several model configurations: a digital reference, a hybrid baseline without noise, a fixed-mismatch noisy baseline, and the proposed Ensemble HAT regime. Additional variants incorporate front-end models and retention-drift corrections for specific diagnostic checks (see Supplementary Note S-Notation).

Nonlinear-write and retention surrogates are anchored to measured DNTT transients \citep{vincze2025dualplasticity}. Nonlinear write modifies the STE backward pass by scaling the straight-through gradient with a state-dependent factor proportional to the present conductance magnitude raised to $NL-1$\footnote{The absence of an explicit $NL$ prefactor is intentional: this is a behavioral proxy for position-dependent update difficulty, not the strict derivative of a power-law $f(G)=G^{NL}$.}; ideal write is recovered when $NL_{\mathrm{LTP}}=1$ (long-term potentiation) and $|NL_{\mathrm{LTD}}|=1$ (long-term depression). The optional optoelectronic front-end model applies inverse-gamma preprocessing, photocurrent shot-noise injection, and per-sample renormalization; the explicit response equations are kept in the Supplementary Information.

\subsection{Profile-Driven Substitution Interface}
\label{subsec:profile-interface}

The profile interface packages technology-specific assumptions into a replaceable JSON parameter bundle. Deployment holds the digital backbone and training recipe fixed while substituting the profile, enabling literature priors today and measured-device calibration later without code changes.

\subsection{Sensitivity and Energy Metrics}
\label{subsec:sensitivity-energy}

For the joint D2D--ADC sweep, first-order factor importance is quantified with Sobol indices estimated from the $7 \times 9$ grid of Monte Carlo means. Energy is estimated via a first-order analytical model (Supplementary Methods) from operator counts and literature-proxy constants ($E_{\text{cell}}=100$~fJ, $E_{\text{ADC},8b}=25$~fJ, $E_{\text{DAC},8b}=30$~fJ); these are analytical proxy values for relative comparison, not measured silicon values.
